% chapters/01_einleitung.tex
% Erstes Kapitel: Einleitung
% Typische Struktur einer Einleitung in einer Bachelorarbeit

\chapter{Einleitung}
\label{chap:einleitung}

% Einstieg: Kontext und Relevanz des Themas (ca. 1 Absatz)
Hier steht der einleitende Text, der das Thema in seinen Kontext einbettet und
die Relevanz der Fragestellung begr\"{u}ndet.

\section{Motivation und Problemstellung}
\label{sec:motivation}

Hier wird das konkrete Problem beschrieben, das die Arbeit adressiert. Was ist
der Ausgangspunkt? Warum ist das ein relevantes Problem?

\section{Zielsetzung der Arbeit}
\label{sec:zielsetzung}

Hier wird das Ziel der Arbeit formuliert. Was soll am Ende der Arbeit
erreicht sein?

\section{Forschungsfragen}
\label{sec:forschungsfragen}

Die Forschungsfragen strukturieren die Untersuchung. Beispiel:

\begin{itemize}
  \item Forschungsfrage 1?
  \item Forschungsfrage 2?
  \item Forschungsfrage 3?
\end{itemize}

\section{Methodisches Vorgehen}
\label{sec:methodik}

Hier wird kurz beschrieben, wie die Forschungsfragen beantwortet werden.
Welche Methoden werden eingesetzt?

\section{Aufbau der Arbeit}
\label{sec:aufbau}

Kapitel~2 behandelt \ldots Kapitel~3 \ldots
